\section{References}

\subsection{Scripture and the liturgical books}

\begin{description}
\item[The Holy Bible.] Douay--Rheims, Challoner revision,
  \href{https://www.gutenberg.org/ebooks/1581}{Project Gutenberg 1581}.
  Mary's canticle in the prayer appendix is Luke 1:46--55.

\item[The Roman Missal.] \textit{Missale Romanum}, Vatican typical edition,
  1962. Sundays I--III, pp.~1--5. Ember Masses, pp.~5--14.
  Sunday IV and the Vigil, pp.~14--16. Masses of the Nativity,
  pp.~17--20. \href{https://musicasacra.com/pdf/missale62.pdf}{Facsimile}.

\item[The general rubrics.] \textit{Rubricae Breviarii et Missalis Romani},
  July~26, 1960, \textit{Acta Apostolicae Sedis} 52, pp.~593--740,
  especially nos.~15, 20, 30, 71--72, 131, 299 and 468.
  \href{https://www.vatican.va/archive/aas/documents/AAS-52-1960-ocr.pdf}{Official text}.

\item[The Roman Breviary.] \textit{The Roman Breviary}, trans. John,
  Marquess of Bute, new edition for use in England, vol.~I, Winter
  (Edinburgh and London: Blackwood, 1908).
  \textit{Alma Redemptoris Mater}, p.~210, the Great O Antiphons,
  pp.~243--244, and the Sunday Matins lessons, pp.~216, 229, 241.
  \href{https://archive.org/details/romanbreviaryref01cathuoft}{Digitized volume}.
\end{description}

\subsection{The Fathers and St Bernard}

\begin{description}
\item[Augustine.] Letter 130 to Proba,
  \S\S18--20, 31, trans. J.~G.~Cunningham, \textit{Nicene and Post-Nicene
  Fathers}, first series, vol.~I (New York, 1886), pp.~465--466, 469.
  \href{https://www.ccel.org/ccel/schaff/npnf101.html}{Volume}.

\item[Bernard of Clairvaux.] Fifth Sermon for Advent, \S\S1--3,
  trans. a Priest of Mount Melleray, \textit{St.~Bernard's Sermons
  for the Seasons}, vol.~I (Dublin: Browne and Nolan, 1921),
  pp.~40--43. Latin in \textit{Patrologia Latina} 183, cols.~50--51.

\item[Cyril of Alexandria.] \textit{Commentary on Luke}, Sermon I,
  Luke 2:7, trans. R.~Payne Smith (Oxford, 1859), p.~12.
  \href{https://www.tertullian.org/fathers/cyril_on_luke_01_sermons_01_11.htm}{Text}.

\item[Cyril of Jerusalem.] \textit{Catechetical Lecture} 15.1--3,
  trans. E.~H.~Gifford, \textit{Nicene and Post-Nicene Fathers},
  second series, vol.~VII (New York, 1894), pp.~104--106.
  \href{https://www.ccel.org/ccel/schaff/npnf207.html}{Volume}.

\item[Gregory of Nazianzus.] \textit{Third Theological Oration}
  (Oration~29), \S19, trans. C.~G.~Browne and J.~E.~Swallow,
  \textit{Nicene and Post-Nicene Fathers}, second series, vol.~VII
  (New York, 1894), p.~308.
  \href{https://www.ccel.org/ccel/schaff/npnf207.html}{Volume}.

\item[Gregory the Great.] \textit{Homilies on the Gospels} 1.1, 3;
  6.1--2; 7.1--3; 8.1; 20.2--4. Latin text in \textit{Patrologia
  Latina} 76 (Garnier reprint, 1878). English excerpts from Homilies
  1, 6 and 7 in Bute's \textit{Roman Breviary} at the pages above.
  Bute's heading for the Second Sunday calls the homily the tenth.
  The passage is from Homily~6. For the Latin, see the
  \href{https://la.wikisource.org/wiki/Homiliarum_in_Evangelia}{complete text}.

\item[Leo the Great.] Sermons 16.2 and 19.1--2 on the December
  fast, and 26.2, 5 on the Nativity, trans. C.~L.~Feltoe,
  \textit{Nicene and Post-Nicene Fathers}, second series, vol.~XII
  (New York, 1895), pp.~124, 127--128, 137--139.
  \href{https://www.ccel.org/ccel/schaff/npnf212.html}{Volume}.
\end{description}

\subsection{The history and discipline of Advent}

\begin{description}
\item[The Roman calendar of 354.] \textit{Depositio Martirum},
  first entry, ed. Theodor Mommsen, \textit{Monumenta Germaniae
  Historica}, \textit{Auctores antiquissimi} 9 (Berlin, 1892),
  pp.~71--72.
  \href{https://www.dmgh.de/mgh_auct_ant_9/index.htm\#page/71/mode/1up}{Edition}.

\item[Councils and bishops.] First Council of Zaragoza, canon~4,
  in Mansi, \textit{Sacrorum conciliorum nova et amplissima collectio},
  vol.~III (Florence, 1759), cols.~633--635.
  Gregory of Tours, \textit{Histories} 10.31, section~VI,
  ed. Wilhelm Arndt (Hanover, 1884), pp.~444--445, and trans. O.~M.~Dalton (Oxford, 1927), vol.~II, p.~471.
  First Council of Mâcon, canon~9, ed. Friedrich Maassen,
  \textit{Concilia aevi Merovingici} (Hanover, 1893), p.~157.

\item[The old Roman Mass prayers.] \textit{The Gregorian
  Sacramentary under Charles the Great}, ed. H.~A.~Wilson
  (London, 1915), pp.~113--116, nos.~150--157.
  \href{https://commons.ptsem.edu/id/gregoriansacrame0049cath}{Digitized volume}.
  Amalarius of Metz, \textit{Liber officialis} 3.40,
  \textit{Patrologia Latina} 105, col.~1219.

\item[Latin Church law.] \textit{Codex Iuris Canonici} (1917),
  can.~1252, \S\S2, 4.
  Paul~VI, \href{https://www.vatican.va/content/paul-vi/en/apost_constitutions/documents/hf_p-vi_apc_19660217_paenitemini.html}{\textit{Paenitemini}}
  (February~17, 1966), norm~II.
  \textit{Code of Canon Law} (1983),
  \href{https://www.vatican.va/archive/cod-iuris-canonici/eng/documents/cic_lib4-cann1244-1253_en.html}{cann.~1249--1253}.

\item[United States.] National Conference of Catholic Bishops,
  \href{https://www.usccb.org/prayer-and-worship/liturgical-year-and-calendar/lent/us-bishops-pastoral-statement-on-penance-and-abstinence}{Pastoral Statement on Penance and Abstinence}
  (November~18, 1966), \S\S4--7, 17, 22--25.
  \href{https://www.usccb.org/beliefs-and-teachings/what-we-believe/canon-law/complementary-norms/canons-1252-and-1253-observance-of-fast-and-abstinence}{Complementary norm to cann.~1252 and 1253}.

\item[Further reading.] Triptych's
  \href{https://mystago.gy/library/faith.html\#the-four-last-things}{Four Last Things}
  studies follow the themes of death, judgment, hell and heaven.
\end{description}
